% R598: whole EGA I Chapter0 3.5.4.
% OpenAI Codex, GPT-6.1 Sol, Ultra; no human review claimed.
\oldpage[0\textsubscript{I}]{32}

\begin{env}[3.5.4]
\label{0.3.5.4-fr}
Supongamos ahora que la categoría $\boldsymbol{K}$ sea tal\footnote{En
el libro citado en la Introducción, daremos condiciones muy
generales sobre la categoría $\boldsymbol{K}$ que aseguren la existencia de las imágenes
inversas de prehaces con valores en $\boldsymbol{K}$.} que \emph{todo}
prehaz $\mathcal{G}$ sobre $Y$ con valores en $\boldsymbol{K}$ admita una
imagen inversa por $\psi$, que denotaremos por $\psi^*(\mathcal{G})$.

Veremos que se puede definir $\psi^*(\mathcal{G})$ como \emph{funtor
covariante} en $\mathcal{G}$, de la categoría de prehaces sobre $Y$ con
valores en $\boldsymbol{K}$ a la de haces sobre $X$ con valores en
$\boldsymbol{K}$, de tal manera que el isomorfismo $v\to v^\flat$ sea un
\emph{isomorfismo de bifuntores}
\begin{equation*}
  \operatorname{Hom}_X(\psi^*(\mathcal{G}),\mathcal{F})
  \xrightarrow{\sim}
  \operatorname{Hom}_Y(\mathcal{G},\psi_*(\mathcal{F}))
\tag{3.5.4.1}\label{0.3.5.4.1-fr}
\end{equation*}
en $\mathcal{G}$ y $\mathcal{F}$.

En efecto, para todo morfismo
$w:\mathcal{G}_1\to\mathcal{G}_2$ de prehaces sobre $Y$ con valores en
$\boldsymbol{K}$, consideremos el morfismo compuesto
\[
  \mathcal{G}_1\xrightarrow{w}\mathcal{G}_2
  \xrightarrow{\rho_{\mathcal{G}_2}}
  \psi_*(\psi^*(\mathcal{G}_2))~;
\]
le corresponde un morfismo
$(\rho_{\mathcal{G}_2}\circ w)^\sharp:
\psi^*(\mathcal{G}_1)\to\psi^*(\mathcal{G}_2)$, que denotaremos por
$\psi^*(w)$. Se tiene, pues, en virtud de (3.5.3.3)
\begin{equation*}
  \psi_*(\psi^*(w))\circ\rho_{\mathcal{G}_1}
  =\rho_{\mathcal{G}_2}\circ w.
\tag{3.5.4.2}\label{0.3.5.4.2-fr}
\end{equation*}
Para todo morfismo $u:\mathcal{G}_2\to\psi_*(\mathcal{F})$, donde
$\mathcal{F}$ es un haz sobre $X$ con valores en $\boldsymbol{K}$, se tiene,
según (3.5.3.3), (3.5.4.2) y la definición de $v^\flat$
\[
  (u^\sharp\circ\psi^*(w))^\flat
  =\psi_*(u^\sharp)\circ\psi_*(\psi^*(w))\circ\rho_{\mathcal{G}_1}
  =\psi_*(u^\sharp)\circ\rho_{\mathcal{G}_2}\circ w
  =u\circ w
\]
o, equivalentemente,
\begin{equation*}
  (u\circ w)^\sharp=u^\sharp\circ\psi^*(w).
\tag{3.5.4.3}\label{0.3.5.4.3-fr}
\end{equation*}
Si se toma, en particular, como $u$ un morfismo
\[
  \mathcal{G}_2\xrightarrow{w'}\mathcal{G}_3
  \xrightarrow{\rho_{\mathcal{G}_3}}
  \psi_*(\psi^*(\mathcal{G}_3)),
\]
resulta
\[
  \psi^*(w'\circ w)
  =(\rho_{\mathcal{G}_3}\circ w'\circ w)^\sharp
  =(\rho_{\mathcal{G}_3}\circ w')^\sharp\circ\psi^*(w)
  =\psi^*(w')\circ\psi^*(w),
\]
de donde se sigue nuestra afirmación.

Finalmente, para todo haz $\mathcal{F}$ sobre $X$ con valores en
$\boldsymbol{K}$, sea $i_{\mathcal{F}}$ el morfismo identidad de
$\psi_*(\mathcal{F})$ y denotemos por
\[
  \sigma_{\mathcal{F}}:
  \psi^*(\psi_*(\mathcal{F}))\to\mathcal{F}
\]
el morfismo $(i_{\mathcal{F}})^\sharp$~; la fórmula (3.5.4.3) da, en
particular, la factorización
\begin{equation*}
  u^\sharp:\psi^*(\mathcal{G})
  \xrightarrow{\psi^*(u)}
  \psi^*(\psi_*(\mathcal{F}))
  \xrightarrow{\sigma_{\mathcal{F}}}\mathcal{F}
\tag{3.5.4.4}\label{0.3.5.4.4-fr}
\end{equation*}
para todo morfismo $u:\mathcal{G}\to\psi_*(\mathcal{F})$. Diremos que el
morfismo $\sigma_{\mathcal{F}}$ es \emph{canónico}.
\end{env}
